\documentclass[11pt]{article}
\usepackage[margin=0.85in]{geometry}
\usepackage{amsmath,amssymb,booktabs,longtable,array,xcolor}
\usepackage[colorlinks,linkcolor=blue!45!black]{hyperref}

\newcommand{\M}{\mathcal M}
\definecolor{blk}{RGB}{155,25,20}
\definecolor{con}{RGB}{150,90,10}
\definecolor{add}{RGB}{20,105,60}
\newcommand{\BLOCK}{\textcolor{blk}{\textbf{BLOCKING}}}
\newcommand{\FIX}{\textcolor{con}{\textbf{FIX}}}
\newcommand{\ADD}{\textcolor{add}{\textbf{ADD}}}

\setlength{\parindent}{0pt}
\setlength{\parskip}{4pt}
\newcommand{\pl}[2]{\textbf{p.#1, l.#2}}

\title{Critical review of \texttt{draft\_close\_to\_final\_aug28}}
\author{}
\date{28 August 2026}

\begin{document}
\maketitle

Page and line numbers refer to the margin numbering of the reviewed PDF. Items are grouped by
what they require: \BLOCK{} must be resolved before submission, \FIX{} is a correction to
existing text, \ADD{} is material that exists but is not in the paper.

Every number quoted below as ``verified'' was recomputed from the saved run outputs.

% ==================================================================
\section{Blocking}

\begin{longtable}{@{}p{0.11\textwidth}p{0.30\textwidth}p{0.53\textwidth}@{}}
\toprule
Location & Issue & Action\\
\midrule
\endhead
\pl{30--36}{960--1305} & The entire NeurIPS checklist is \texttt{[TODO]}: all 16 answers and
all 16 justifications. The instruction block that the template says to delete is still present
(\pl{30}{960--988}). & The checklist itself states that papers without it are desk rejected. Fill
all 16, delete the instruction block. A complete draft with section pointers is available.\\
\addlinespace
\pl{24--29}{769--959} & Appendix C duplicates Appendix B in full: threshold, reference ladder,
guidance, capacity, CLIP and super-resolution each appear twice, with \emph{different numbers}.
The note at \pl{24}{770--773} acknowledges this. & Merge into B and delete C. Beyond reviewer
confusion, iThenticate compares the document against itself; internal duplication inflates the
similarity score directly.\\
\addlinespace
\pl{12}{427--435}, \pl{14}{498--502}, \pl{19}{610--616}, \pl{22}{711--714}, \pl{22}{735--739},
\pl{23}{754--757}, \pl{23}{765--768}, \pl{28}{927--929}, \pl{29}{942--944} &
Nine \texttt{[TODO---Yassine]} blocks remain in the appendices. & All nine are answered in the
material already prepared. None requires new computation.\\
\addlinespace
\pl{22}{722} & A literal \texttt{:contentReference}\discretionary{}{}{}\texttt{[\ldots]}%
\discretionary{}{}{}\texttt{index=0} string appears in the body text of \S B.6, at the end of the
paragraph beginning ``Relative to the unguided model''. & Delete. This is a paste artefact.\\
\bottomrule
\end{longtable}

% ==================================================================
\section{The paper contradicts itself}

Each row is the same quantity reported with two different values in the same document.

\begin{longtable}{@{}p{0.20\textwidth}p{0.24\textwidth}p{0.24\textwidth}p{0.26\textwidth}@{}}
\toprule
Quantity & States & Also states & Correct value\\
\midrule
\endhead
Guidance gain, $w=3$ &
$+0.0770$ / $-0.4905$ \pl{7}{234}, Table C6 \pl{27}{} &
$+0.157$ / $-0.505$ \pl{22}{719--721} &
$+0.0770_{\pm0.0004}$ / $-0.4905_{\pm0.0039}$\\
\addlinespace
Auxiliary-noise sweep &
$0.2149\!\to\!0.3105$, Table B5 \pl{20}{} &
$0.2195\!\to\!0.3122$, $\M_1=0.3195$ \pl{27}{881--882} &
Table B5 is correct\\
\addlinespace
Integration $K\!=\!1\!\to\!64$ &
$0.3541\!\to\!0.3172$, Table B6 \pl{21}{} &
$0.350\!\to\!0.319$ \pl{7}{242--243}, \pl{24}{790--791} &
Table B6 is correct\\
\addlinespace
Capacity gap &
$0.0892$ / $0.1018$ at $\eta{=}10^{-3}$, Table B2 \pl{17}{} &
$0.0985$ / $0.1040$, Table B3 \pl{18}{}, \pl{7}{227}, \pl{26}{851} &
Different sweeps; see \S3\\
\addlinespace
$\M_1$ terminal fidelity &
$0.3184$ \pl{7}{217} &
$0.322$ \pl{6}{209}, $0.3138$ Table B3 &
See the noise floor below\\
\bottomrule
\end{longtable}

\textbf{Evaluation-noise floor, and why several of these are not errors.} The evaluation routine
never fixes its own seed, while the distributional routine does. Two rows of the saved output ---
the $\varepsilon=0$ row of Table B4 and the $\sigma_{\mathrm{aux}}=0$ row of Table B5 --- are the
\emph{same trained network} (retraining at $\sigma_{\mathrm{aux}}=0$ with the default seed
reproduces it, as their byte-identical $W_1=0.036684$ and energy distance $0.537162$ confirm).
Their cosines nevertheless differ by $0.0016$. \textbf{Differences below roughly $0.002$ in any
reported cosine are evaluation noise.} Stating this once, in the evaluation protocol, resolves the
$\M_1$ spread and several smaller discrepancies without further edits. Every gap the paper claims
is between $0.04$ and $0.12$, forty to seventy times larger.

\textbf{Capacity: the two tables are different experiments.} Table B2 (\pl{17}{}) is the
three-learning-rate grid, $120$ trainings. Table B3 (\pl{18}{}) is the earlier width sweep at a
single automatically-selected learning rate. They are not reconcilable as stated because they are
not the same sweep. Recommendation: keep B2, delete B3, and move the training-loss and
distributional-error trend (\pl{18}{598--603}) into B2's discussion. The main text at
\pl{7}{226--229} then quotes B2 rather than B3.

% ==================================================================
\section{Numbers superseded by verified runs}

\begin{longtable}{@{}p{0.20\textwidth}p{0.30\textwidth}p{0.44\textwidth}@{}}
\toprule
Location & Currently & Verified\\
\midrule
\endhead
\pl{7}{252--253}, \pl{22}{746--748}, Table B7 \pl{23}{}, \pl{28}{894--895} &
CLIP alignment $0.50\to0.69$, prototype $0.79$ &
$0.5517\to0.7586$, prototype $0.7930$, random-pair floor $0.4815$. Two independent runs (the
sweep and the probe rerun) agree.\\
\addlinespace
\pl{22}{750} & ``the measured factorization discrepancy is $0.62$'' &
No quantity of that name is computed. The stage was re-run to check: no column equals $0.62$.
Closest is $\hat\beta=0.6010$, already named. \textbf{Delete the sentence.}\\
\addlinespace
\pl{7}{264}, \pl{28}{922} & ``At both completed resolutions'' &
All three are complete. At $96\times96$: PSNR $23.1984$ vs $23.1984$ --- \emph{identical}, not
``slightly higher'' --- energy $1.888$ vs $2.222$, $W_1^{\mathrm{hf}}$ $623.6$ vs $659.4$.\\
\addlinespace
\pl{16}{571--573}, \pl{25}{810--812} & ``maximum absolute error of $3.3\times10^{-2}$'' &
Correct, but the error is \emph{one-sided}: measured exceeds analytic in $22$ of $27$ cells, mean
signed error $+0.0099$, sign test $p=1.5\times10^{-3}$.\\
\bottomrule
\end{longtable}

The $96\times96$ row is the strongest single result in the super-resolution section and is
currently absent: the two models are indistinguishable on the pointwise metric to four decimals
while the distributional gap persists at $18\%$. Add it to Table C8/B8.

% ==================================================================
\section{Experiments that were run and are not in the paper}

Six completed experiments do not appear anywhere in the draft. Two of them change what the paper
can claim.

\subsection*{4.1 The criterion verified at scale (job 381574)}

\ADD{} --- belongs in \S A.2.1, after \pl{11}{388--391}.

The paper currently supports Proposition~1 with a $180$-point power-law fit. A separate run
tested the criterion directly across dimensions $n\in\{2,4,8,16,32,64\}$:

\begin{center}
\begin{tabular}{@{}lr@{}}
\toprule
Systems satisfying both hypotheses, worst $90$th-percentile drift & $4.06\times10^{-8}$\\
Systems violating a hypothesis at relative size $0.1$, smallest median drift & $9.78\times10^{-3}$\\
\textbf{Separation} & $\mathbf{2.41\times10^{5}}$\\
\bottomrule
\end{tabular}
\end{center}

This is a far stronger statement than the power-law fit alone and costs two sentences.

\subsection*{4.2 The two hypotheses separated numerically}

\ADD{} --- belongs immediately after Proposition~1, \pl{3}{119--123}.

Theorem~4.11 has two hypotheses. A four-arm experiment isolates each:

\begin{center}
\begin{tabular}{@{}lrl@{}}
\toprule
Arm & max relative drift & \\
\midrule
symmetric, in span, commuting marginals    & $6.34\times10^{-5}$ & invariant\\
symmetric, in span, non-commuting marginals& $6.34\times10^{-5}$ & invariant\\
symmetric, off span                        & $7.62\times10^{-3}$ & moves, $120\times$\\
skew, in span                              & $9.52\times10^{-1}$ & moves, $15{,}000\times$\\
skew, off span                             & $1.01$              & moves\\
\bottomrule
\end{tabular}
\end{center}

Two things the paper asserts but does not evidence follow directly. First, the Remark at
\pl{11}{381--383} --- that invariance does not require $[\Sigma_0,\Sigma_1]=0$ --- is confirmed
to five digits by the first two rows. Second, the skew/span asymmetry that \S A.2.1 quantifies as
$6.585$ is visible here as a $125\times$ ratio between the two single-violation arms.

\subsection*{4.3 A masked-diffusion analogue --- the effect does not transfer}

\ADD{} --- belongs in \S8 Limitations, \pl{8}{297--305}, and as a short appendix subsection.

This is the most consequential omission. A discrete masked-diffusion analogue was run over six
source-informativeness levels, with four arms crossing ``start from mask or from source'' with
``label only or sees source'':

\begin{center}
\small
\begin{tabular}{@{}rrrrrr@{}}
\toprule
$\beta_d$ & do nothing & mask start, & mask start, & source start, & source start,\\
          &            & label only  & sees source & label only    & sees source\\
\midrule
$0.10$ & $0.612$ & $0.806$ & $0.854$ & $0.852$ & $0.850$\\
$0.25$ & $0.612$ & $0.713$ & $0.820$ & $0.826$ & $0.830$\\
$0.40$ & $0.612$ & $0.606$ & $0.819$ & $0.809$ & $0.820$\\
$0.55$ & $0.613$ & $0.458$ & $0.817$ & $0.803$ & $0.820$\\
$0.70$ & $0.613$ & $0.355$ & $0.831$ & $0.801$ & $0.821$\\
$0.85$ & $0.611$ & $0.211$ & $0.844$ & $0.799$ & $0.812$\\
\bottomrule
\end{tabular}
\end{center}

The comparison that matters is the last two columns: conditioning on the source while starting
from it. In the continuous spherical model this costs roughly $0.10$ in terminal fidelity. Here
it \emph{helps}, by $+0.0107$ on average and at five of six levels.

\textbf{The source-conditioning penalty is therefore specific to the continuous deterministic
setting and does not transfer to a discrete masked model.} The paper's limitation paragraph
(\pl{8}{303--305}) currently says only that whether the mechanism persists ``under alternative
paths, stochastic dynamics, and high-dimensional image-space flow matching remains open''. It is
not open --- it was tested, and in one such setting the sign reverses. Reporting this is
substantially better than having a reviewer ask.

\subsection*{4.4 The parameterisation control}

\ADD{} --- belongs in \S B.2, near \pl{17}{583--585}.

The capacity conclusion currently rests on maximal-update parameterisation alone. The standard
parameterisation was also run: at width $2048$ the gap is $0.0998$ under standard and $0.0988$
under maximal-update. The conclusion is therefore not a parameterisation artefact, which is
exactly what a reviewer will ask.

\subsection*{4.5 The integrator anchor test}

\ADD{} --- belongs in \S B.3, near \pl{20}{658--664}.

The $\alpha$-family and the auxiliary-noise family were originally sampled with different
integrators. At $\alpha=0$ and $\sigma_{\mathrm{aux}}=0$ both branches use the same trained model
(terminal losses $4.9116\times10^{-4}$ and $4.9046\times10^{-4}$), yet reported distributional
errors of $0.0782$ and $0.0367$. Re-run under one integrator, the anchor agrees and the two
families trace one curve (cross/within ratio $1.40$ matched against $3.44$ mismatched). If the
paper reports both families as one comparison, it needs this sentence.

\subsection*{4.6 Six-factor super-resolution sweep, and the Stable Diffusion rescore}

\ADD{} --- \S B.8 and the transfer paragraph at \pl{28}{930--941}.

\emph{Factor sweep.} At fixed resolution $48$, six downsampling factors were run. The source-rank
bound is hit exactly at every one ($1728, 768, 432, 192, 108, 48$), taking the total to
\textbf{nine exact agreements}. The trained flow beats doing nothing at every factor, with the
gain \emph{falling} as the source degrades ($3.81, 2.26, 1.97, 1.29, 1.02, 0.59$ PSNR).

That is the opposite direction from the pretrained model at \pl{28}{930--941}, and the contrast
is the useful part: \textbf{a flow trained for its exact degradation beats doing nothing
everywhere tested; a pretrained flow applied off the shelf is worse than doing nothing at
$2\times$, $4\times$ and $8\times$}. The do-nothing regime lives in the mismatch, not in the task.

\emph{Distributional rescore.} \S8 states that single-target cosine ``should not be interpreted
as a standalone measure of distributional fidelity'' (\pl{8}{302--303}), yet the Stable Diffusion
experiment is scored by CLIP cosine alone. It has now been rescored with four samples per image:

\begin{center}
\small
\begin{tabular}{@{}rrrrr@{}}
\toprule
degradation & source align & best cosine gain & best energy gain & pairwise cosine at $s=0.8$\\
\midrule
$2\times$  & $0.997$ & $-0.189$ & $-0.174$ & $0.742$\\
$4\times$  & $0.853$ & $-0.061$ & $+0.008$ & $0.736$\\
$8\times$  & $0.733$ & $-0.055$ & $+0.054$ & $0.719$\\
$16\times$ & $0.430$ & $+0.120$ & $+0.392$ & $0.706$\\
\bottomrule
\end{tabular}
\end{center}

Null floor for the energy distance is $0.011$; the retained source has pairwise cosine $1.000$ by
construction. Three consequences. At $2\times$ the do-nothing regime holds on \emph{both}
instruments, so the headline claim no longer depends on the discredited metric. At $4\times$ the
energy gain of $+0.008$ is below the null floor and should not be read as a reversal. At
$8\times$ the two instruments genuinely invert --- cosine $-0.055$, energy $+0.054$ --- which is
a \textbf{fourth instance of the paper's own mean-seeking inversion, in a production model}.

\subsection*{4.7 The held-out label probe}

\ADD{} --- \S B.7, near \pl{27}{885--893}.

The probe reaches $0.058$ held-out accuracy against a chance level of $0.10$, on a disjoint
$80/20$ split. Phrase it as ``does not extract label information'', not as an accuracy figure:
the value sits about five standard errors \emph{below} chance, which indicates the probe
anti-generalises rather than that it is informative.

% ==================================================================
\section{Mathematics available but not included}

\subsection*{5.1 The threshold criterion is stronger than stated}

\FIX{} --- \pl{4}{153--157}, \pl{14}{494--502}, Table B1 \pl{16}{}, \pl{16}{571}.

The TODO at \pl{14}{498--502} asks whether $\Delta$ can have a root when $A\ge\cos\jmath$. It
always does. Substituting $\beta=\sin\varphi$ and $\psi(\varphi)=\arccos(A\cos\varphi)$,
\[
\Delta/A=\cos(\psi-\varphi)-\cos(\varphi-\jmath),
\]
verified to $7\times10^{-15}$. At $\varphi=\jmath$ the second term is $\cos 0=1$, its maximum,
while $\psi(\jmath)>\jmath$ strictly for $A<1$; hence $\Delta(\sin\jmath)<0$ \emph{for every}
$A\in(0,1)$, and with $\Delta(1)>0$ a transition above $\sin\jmath$ exists unconditionally.
Uniqueness is what the criterion governs: $\psi'\le A$ makes
$h=2\varphi-\jmath-\psi$ strictly increasing, giving exactly one root above $\sin\jmath$; below
it the condition reduces to $A\cos\varphi>\cos\jmath$, monotone, so a second root appears exactly
when $A>\cos\jmath$. Confirmed on $120$ cells with no exceptions.

Consequences: Theorem~2 should state existence unconditionally and $A<\cos\jmath$ as the
uniqueness condition; the three dashes in Table B1 should be replaced by the measured upper roots
$0.8286$, $0.5600$, $0.7586$, which agree with the analytic values to $\le 0.027$; and the
comparison becomes \textbf{30 cells rather than 27}.

\subsection*{5.2 The first-order derivation, and the asymptotic slope}

\FIX{} --- \pl{12}{427--435}.

The full derivation exists: the identity $G_0=\tfrac12V_0'$; the homogeneous-marginal limit in
which every symmetric response vanishes exactly by parity and the skew response has the closed
form $-I_\kappa N\Sigma^{-1}$ with
$I_\kappa=\tfrac{4}{\sqrt{\mu(4-\mu)}}\arctan\sqrt{\mu/(4-\mu)}$; and the pairwise commuting model
giving the suppression bound
$\mathcal B_{ij}/|\mathcal A_{ij}|\ge 8/\sup_t|(\log v_i/v_j)'|$.

On the asymptotic slope at \pl{12}{426}: the ratio slope is identically the skew slope minus the
span slope, an identity holding in the data to $10^{-15}$. Over the reported window the span
exponent is $0.573$ and the skew exponent $-0.159$, giving $-0.732$; a decade lower they are
$0.933$ and $-0.019$, giving $-0.952$. So $-1$ is asymptotic and the shortfall is the span
response being sublinear over a pre-asymptotic window --- \emph{not} the skew constant's decline,
which pushes toward $-1$.

\subsection*{5.3 The $1/t$ mechanism on the geodesic}

\ADD{} --- \pl{15}{521--545}, \pl{8}{303--305}.

Propositions~2 and 3 and \S A.6 are derived for the linear interpolant
$z_t=(1-t)z_0+tz_1$, while every experiment uses the geodesic of (3.4). \S8 flags this as open.
It is not: with $z_t$ and the auxiliary state both in $\mathrm{span}\{z_0,z_1\}$, the recovery
determinant collapses to $-\sin((1-\alpha)t\theta)/\sin\theta$ and the sensitivity is
\[
\frac{\partial z_1}{\partial z_t}=\frac{\sin((1-\alpha t)\theta)}{\sin((1-\alpha)t\theta)},
\]
which diverges as $1/t$ exactly as in the flat case, differing only by the bounded factor
$\sin\theta/\theta\in(0,1]$; the remainder is second order in $t$ for $\alpha=0$. Verified against
finite differences in ambient dimension $64$ to $2.7\times10^{-10}$, and the truncation
experiment's $\varepsilon$-ordering is identical under both interpolants. This closes the gap
between \S6 and the retraction of flat-to-sphere transfer elsewhere in the paper.

% ==================================================================
\section{Figures}

\begin{longtable}{@{}p{0.13\textwidth}p{0.36\textwidth}p{0.45\textwidth}@{}}
\toprule
Location & Issue & Action\\
\midrule
\endhead
Figure 1, \pl{6}{206} & Three panels at one-third column width; axis labels and legends are
unreadable at print size. & Use the three panels as separate figures, or give the composite a
full page width.\\
\addlinespace
Figure 2, \pl{17}{574} & Plots $j=0$ points, which Table B1's own caption calls ``formal only'';
also plots the three $A\ge\cos\jmath$ cells that have no analytic counterpart, which is why
points appear at $A=0.9$ with no curve. & Restrict to the $27$ cells with a predicted threshold,
or --- better --- adopt \S5.1 above and plot all $30$ with their upper roots.\\
\addlinespace
Figure 3, \pl{18}{591} & Legend entries are malformed: ``source-conditioned $\eta$:0.000 1''. &
Regenerate.\\
\addlinespace
Figures 1 and 5 & Identical figure appears twice (\pl{6}{206} and \pl{24}{788}). & Consequence of
the duplicated appendix.\\
\addlinespace
--- & \textbf{The paper contains no qualitative figures} despite two image experiments. &
Two panels are available: super-resolution (degraded / regression / $\M_1$ / $\M_2$ / target) in
which $\M_2$ is visibly smoother than $\M_1$ while scoring higher, and the regression arm
sharpest of all while scoring worst; and Stable Diffusion (clean / retained / three strengths per
degradation). The first makes the central metric argument visible rather than asserted.\\
\bottomrule
\end{longtable}

% ==================================================================
\section{Smaller items}

\begin{longtable}{@{}p{0.15\textwidth}p{0.80\textwidth}@{}}
\toprule
Location & Comment\\
\midrule
\endhead
\pl{6}{209}, Table C1 & ``$0.137$ (perfect sampler)''. Two distinct references exist:
$A^2=0.1369$ for a fresh draw sharing the posterior mean, and $(As)^2=0.0493$ for a fresh draw
given $c$ alone. The measured true-posterior row is $0.0500$. State which is meant; reporting
both is clearer, since the gap between them is itself informative about what the metric rewards.\\
\addlinespace
\pl{9}{306--332} & Nine references. Thin for the venue, and an Extended Related Work is planned.
The prior-art audit already covers sixteen items with quoted statements and identified
differences; condensing it is not new research.\\
\addlinespace
--- & Nine sources in the prior-art record are marked as unverified secondary summaries, with a
standing note that each must be confirmed before submission. The load-bearing one is the prior
for the $1/t$ mechanism in \S6/\S A.6. If it proves more than the summary suggests, that
section's contribution narrows. This is an hour of reading, not an experiment.\\
\addlinespace
\pl{8}{303--305} & The limitation ``whether the same mechanism persists under \dots\ stochastic
dynamics \dots\ remains open'' should be revised in light of \S4.3 above.\\
\addlinespace
\pl{18}{604--607}, \pl{26}{853--855} & The $4\times$ training result ($0.215\to0.161$) comes from
the earlier single-learning-rate sweep. State which learning rate, or re-run it inside the
three-rate grid.\\
\bottomrule
\end{longtable}

% ==================================================================
\section{Suggested order of work}

\begin{enumerate}
\item Merge Appendix C into B and delete C. This removes most of the self-contradictions in
\S2 automatically, and is the single largest reduction in similarity-check exposure.
\item Fill the nine TODOs from the prepared material, and delete the paste artefact at
\pl{22}{722}.
\item Apply the superseded numbers in \S3, and add one sentence on the evaluation-noise floor to
the evaluation protocol.
\item Adopt the strengthened threshold statement (\S5.1); it fills three table cells and takes
the comparison to $30$.
\item Add the two figures and the masked-diffusion limitation (\S4.3) --- the two additions most
likely to be raised by a reviewer if absent.
\item Complete the checklist.
\end{enumerate}

Items 1--3 are corrections to existing text. Items 4--5 add material that is already computed.
Nothing in this document requires a new experiment.

\end{document}
